\documentclass[11pt]{article}
\usepackage[margin=1in]{geometry}
\usepackage{url}
\usepackage{parskip}
\setlength{\parindent}{0pt}

\begin{document}
\thispagestyle{empty}

Dear Editors,

Please find enclosed our submission entitled

\begin{center}
\textbf{``Boyd's conductor-11 Mahler measure conjecture: proof of the
split-integral identity (C3), with an exact structural analysis of the
family $S_k$''}
\end{center}

by Huimin Zheng (College of Information and Network Engineering, Anhui
Science and Technology University, Fengyang, Anhui 233100, P.~R.~China;
\url{zhhm@ahstu.edu.cn}), which we submit for consideration in
\emph{The Ramanujan Journal}.

This manuscript is original, has not been published previously, and is
not under consideration by any other journal. All authors (a single
author) are aware of and approve this submission.

\textbf{Summary and significance.} Boyd's 1998 tables of conjectural
identities between Mahler measures of two-variable polynomials and
$L$-values of elliptic curves begin with the smallest possible
conductor, $N=11$. Of the three conductor-11 identities, the first two
were proved by Brunault; the third, (C3), concerns a polynomial
vanishing on the unit torus and asserts that a signed split integral of
$\log|y|$ around the branch cut equals $b_{11}=L'(E_{11},0)$. This
paper proves (C3), completing the conductor-11 block of Boyd's table.
The main contributions are:
\begin{itemize}
\item A closed-cycle construction: Samart's signed open chain is closed
by a small-branch compensating arc into an anti-invariant integral
cycle, whose homology class is proved to be twice a generator of the
anti-invariant homology; the period ratio is a-priori an integer and
is pinned to $2$ by certified interval arithmetic (Arb ball arithmetic
with rigorous error bounds).
\item A direct regulator computation via Brunault's proved Siegel-unit
formula --- the symbol $\{x,y\}$ is a pair of modular units on
$X_1(11)$, so Bloch's diamond theorem (and its normalization
ambiguities) is entirely bypassed.
\item A complete structural analysis of the family $S_k$: exact
determination of the torus intersections, the unique modular-unit
case, and the tempered cases, together with a mechanistic
counterexample at conductor $53$ showing precisely where the
closed-cycle/modular-unit mechanism provably fails.
\end{itemize}

The paper's blend of explicit arithmetic geometry, special values of
$L$-functions, and rigorously certified computation sits squarely
within the scope of \emph{The Ramanujan Journal}.

\textbf{Transparency.} The research was carried out by the author with
the assistance of the AI system Kimi (Moonshot AI); a formal
Declaration on the use of AI tools appears on the title page of the
article, and the author has verified all mathematical content and
takes full responsibility for it. All certification scripts are openly
available at \url{https://github.com/huiminZheng-collab/boyd-conductor11}
and archived at \url{https://doi.org/10.5281/zenodo.21820650}. A
preprint has been submitted to arXiv (identifier to be announced;
available on request).

\textbf{Suggested referees.}
\begin{itemize}
\item Matilde N. Lal\'in (Universit\'e de Montr\'eal)
\item Fran\c{c}ois Brunault (ENS de Lyon)
\item Mathew Rogers
\item Wadim Zudilin (Radboud University)
\item Detchat Samart
\item Marie-Jos\'e Bertin
\end{itemize}
The author has no conflicts of interest with the suggested referees.

Thank you for your consideration.

\vspace{2em}
Sincerely,

Huimin Zheng\\
Anhui Science and Technology University\\
\url{zhhm@ahstu.edu.cn}

\end{document}
